\section{Results}

All five sub-hypotheses pass. The core claim is validated.

\subsection{Main Results}

\textbf{Existence.} The probe achieves \textbf{AUROC = 0.885} on full-scale evaluation (9,500 train / 1,700 val samples), far exceeding the 0.60 threshold.

\textbf{Layer sweep.} We conduct a preliminary layer sweep on a reduced sample (500 train / 200 val) to identify optimal extraction depth. The inverted-U pattern is confirmed:

\begin{table}[h]
\centering
\caption{Layer sweep results showing inverted-U pattern with peak at L15 (50\% depth).}
\label{tab:layer_sweep}
\begin{tabular}{@{}lccc@{}}
\toprule
Layer & Depth (\%) & AUROC & 95\% CI \\
\midrule
L3    & 12.5      & 0.629 & [0.53, 0.73] \\
L7    & 25.0      & 0.736 & [0.65, 0.82] \\
L11   & 37.5      & 0.818 & [0.74, 0.88] \\
\textbf{L15}   & \textbf{50.0}      & \textbf{0.852} & [0.79, 0.90] \\
L18   & 60.0      & 0.822 & [0.75, 0.88] \\
L23   & 75.0      & 0.785 & [0.69, 0.86] \\
L27   & 87.5      & 0.758 & [0.67, 0.83] \\
L31   & 100.0     & 0.766 & [0.67, 0.85] \\
\bottomrule
\end{tabular}
\end{table}

Peak at L15 (50\% depth). Middle $>$ Final: L18 (0.822) $>$ L31 (0.766). The final probe trained on full data achieves 0.885 AUROC, improving over the layer sweep due to increased training data.

\textbf{Baseline comparison.} The probe substantially outperforms output-level metrics:

\begin{table}[h]
\centering
\caption{Comparison of hidden state probe against output-level baselines.}
\label{tab:baseline_comparison}
\begin{tabular}{@{}lcc@{}}
\toprule
Method & AUROC & $\Delta$ vs Probe \\
\midrule
\textbf{Probe (L15)} & \textbf{0.885} & --- \\
Token Entropy & 0.623 & -0.262 \\
Sequence NLL & 0.589 & -0.296 \\
\bottomrule
\end{tabular}
\end{table}

+26.2 AUROC points over token entropy. Hidden states encode correctness information absent from output distributions.
